\chapter*{Appendices}
\addcontentsline{toc}{chapter}{Appendices}

\section*{Appendix A: Serial Communication Protocol Summary}
\setcounter{table}{0}
\renewcommand{\thetable}{A.\arabic{table}}

The host application communicates with the STM32 firmware over a compact binary UART protocol. The purpose of this appendix is to summarize the protocol surface used by the project without moving low-level implementation details into the main chapters.

Each packet starts with a synchronization byte, followed by a command or response code, payload length, optional payload bytes, and a checksum byte. The checksum is calculated as the XOR of the preceding packet bytes. The maximum payload size is 64 bytes, and the vault key payload is 32 bytes.

\begin{table}[!htbp]
\centering
\caption{\label{tab:appendix-uart-packet}UART packet format summary.}
\begin{tabularx}{\textwidth}{lXX}
\toprule
Field & Size & Description\\
\midrule
Sync & 1 byte & Fixed packet start marker, \texttt{0xAA}\\
Code & 1 byte & Command, event, or response code\\
Length & 1 byte & Number of payload bytes\\
Payload & 0--64 bytes & Command or response data, depending on the packet code\\
Checksum & 1 byte & XOR checksum over sync, code, length, and payload bytes\\
\bottomrule
\end{tabularx}
\end{table}

\begin{table}[!htbp]
\centering
\caption{\label{tab:appendix-uart-commands}Host-to-firmware command summary.}
\begin{tabularx}{\textwidth}{lXX}
\toprule
Command & Code & Purpose\\
\midrule
Enroll start & \texttt{0x01} & Starts fingerprint enrollment for the requested slot or automatic slot selection\\
Verify start & \texttt{0x02} & Starts fingerprint verification and returns the matched authorized fingerprint if successful\\
Delete & \texttt{0x03} & Deletes an enrolled fingerprint template from the sensor slot\\
Cancel & \texttt{0x04} & Cancels an active fingerprint operation\\
Get key & \texttt{0x05} & Requests vault key material after successful verification\\
List enrolled & \texttt{0x06} & Lists enrolled fingerprint slots known by the firmware/sensor flow\\
Status & \texttt{0x07} & Checks firmware/sensor availability and status\\
\bottomrule
\end{tabularx}
\end{table}

\begin{table}[!htbp]
\centering
\caption{\label{tab:appendix-uart-responses}Firmware event and response summary.}
\small
\begin{tabularx}{\textwidth}{lXX}
\toprule
Code group & Codes & Meaning\\
\midrule
User interaction events & \texttt{0x10--0x14} & Waiting for finger, finger detected, remove finger, place again, or processing\\
Successful result & \texttt{0x20} & Operation completed successfully\\
Expected failure results & \texttt{0x21--0x28} & No finger, bad image, no match, already exists, slot full, invalid ID, cancelled, or timeout\\
Key data & \texttt{0x29} & Returns 32 bytes of key material after authorization\\
Busy or hardware error & \texttt{0x2A--0x2B} & Firmware is busy or the sensor reported an error\\
\bottomrule
\end{tabularx}
\end{table}

\section*{Appendix B: Storage Tool Screenshots}
\setcounter{figure}{0}
\renewcommand{\thefigure}{B.\arabic{figure}}

The main testing chapter includes command-line evidence from \texttt{lsblk} and \texttt{cryptsetup luksDump}. The following GParted screenshots are included as additional visual evidence of the physical USB storage layout used during testing.

\begin{figure}[!htbp]
    \centering
    \includegraphics[width=0.98\textwidth,height=0.72\textheight,keepaspectratio]{gparted.png}
    \caption{\label{fig:gparted-storage-view-one}GParted view of the physical USB storage after vault formatting.}
\end{figure}

\begin{figure}[!htbp]
    \centering
    \includegraphics[width=0.98\textwidth,height=0.72\textheight,keepaspectratio]{gparted2.png}
    \caption{\label{fig:gparted-storage-view-two}Additional GParted view showing the formatted storage layout.}
\end{figure}
